\hypersetup{%
	pdftitle = {Modelling single cell dynamics with trajectories and gene regulatory networks},
	pdfsubject = {A subject},
	pdfkeywords = {import keywords},
	pdfauthor = {Robrecht Cannoodt},
	pdfcreator = {\LaTeX},
}

%%%%%%%%%%%%%%%%%%%%%%%%
%%% DEFINE TITLEPAGE %%%
%%%%%%%%%%%%%%%%%%%%%%%%
\begin{titlepage}
	
	\includegraphics[height=1.5cm]{fig/logos/logo_UGent_EN_RGB_color}
	
	\centering
	
	\vfill
	{\Huge \textsf{Modelling single cell dynamics\\with trajectories and gene regulatory networks}}
	
	\vspace{2.25cm}
	
	{\LARGE \textsf{Robrecht Cannoodt}}
	
	\vspace{2.25cm}
	
	{\textsf{Thesis submitted in fulfilment of the requirements for the degree of \\Doctor in Computer Science, 2019}}
	
	\vfill
	
	{\textsf{Supervisors: \\
			Prof. Dr. Yvan Saeys\\Prof. Dr. Katleen De Preter}}
	
	\vfill
	
	\includegraphics[height=.7cm]{fig/logos/icon_UGent_WE_EN_RGB_color}\hfill
	\includegraphics[height=1cm]{fig/logos/vib_rf_inflammation_research_rgb_pos} \hfill
	\includegraphics[height=.8cm]{fig/logos/dambi_logo} \hfill
	\includegraphics[height=.7cm]{fig/logos/cmgg_myshort} \hfill
	\includegraphics[height=.7cm]{fig/logos/FWO_Logo_KleurCMYK} 
	\vspace{1cm}
	
	{\textsf{Vakgroep Toegepaste Wiskunde, Informatica, en Statistiek\\
			Faculteit Wetenschappen, Universiteit Gent \\
			Krijgslaan 281 - S2, 9000 Gent}}
		
\end{titlepage}

%%%%%%%%%%%%%%%%%%%%%%%%
%%% ACKNOWLEDGEMENTS %%%
%%%%%%%%%%%%%%%%%%%%%%%%

%\newpage{\thispagestyle{empty}\cleardoublepage}
%\chapter*{Acknowledgements}

\newpage{\thispagestyle{empty}\cleardoublepage}

%%%%%%%%%%%%%
%%% LISTS %%%
%%%%%%%%%%%%%
\newpage{\thispagestyle{empty}\cleardoublepage}
\tableofcontents 

%\newpage{\thispagestyle{empty}\cleardoublepage}
%\listoffigures

%\newpage{\thispagestyle{empty}\cleardoublepage}
%\listoftables

\newpage{\thispagestyle{empty}\cleardoublepage}


%%%%%%%%%%%%%%%%%%%%
%%% NOMENCLATURE %%%
%%%%%%%%%%%%%%%%%%%%
\makenomenclature
\printnomenclature

\nomenclature{HCA}{Human Cell Atlas}
\nomenclature{TCGA}{The Cancer Genome Atlas}

\nomenclature{DNA}{Deoxyribonucleic Acid}
\nomenclature{RNA}{Ribonucleic Acid}
\nomenclature{mRNA}{Messenger RNA}
\nomenclature{TF}{Transcription Factor}

\nomenclature{GRN}{Gene Regulatory Network}
\nomenclature{NI}{Network Inference}
\nomenclature{CART}{Classification And Regression Trees}
\nomenclature{RF}{Random Forests}
\nomenclature{TI}{Trajectory Inference}

\nomenclature{$k$NN}{$k$-nearest-neighbour}
\nomenclature{DR}{Dimensionality Reduction}
\nomenclature{PCA}{Principal Component Analysis}
\nomenclature{MDS}{Multi-Dimensional Scaling}
\nomenclature{t-SNE}{t-distributed Stochastic Neighbor Embedding}
\nomenclature{UMAP}{Uniform Manifold Approximation and Projection}
\nomenclature{LDMS}{Landmark Multi-Dimensional Scaling}
\nomenclature{DE}{Differential Expression}
\nomenclature{TDE}{Trajectory Differential Expression}

\chapter*{Summary}

Recent technological developments triggered a major milestone in cellular biology by being able to profile and monitor biomolecular contents of single cells at a genome-wide level. The rapid advancements in single-cell omics allow for new types of analyses that fundamentally redefine our understanding of cellular heterogeneity and dynamic processes. Even now, there are even more exciting advancements on the horizon, with the Human Cell Atlas aiming to create a biomolecular and spatial map of all the different cell types in the body. 

Throughout this doctoral project, I did not only get to follow these inspiring advancements first hand, but was able to make scientific contributions to the field by extensively reviewing novel computational tools, and developing several tools and algorithms of my own. More specifically, my efforts have resulted in the following contributions.

A common problem of pioneering computational tools is that during their development, there are rarely sufficient datasets available for adequately quantitatively assessing its performance. We developed a simulator of \textit{in silico} single cells aimed at being flexible and extendible enough to be used to help kick-start emerging computational sub-domains (Chapter~\ref{chap:dyngen}).

We performed a comprehensive comparison of 45 trajectory inference methods (Chapter~\ref{chap:dynbenchmark}). Trajectory inference is a new type of computational analysis whose aim is to identify and study transitions between main cellular states. While trajectory inference is one of the most rapidly growing new fields in single-cell omics tools, a quantitative study of the advantages and disadvantages was hitherto lacking.

We developed four new software tools ourselves. These methods include an extensive toolkit for inferring, visualising and interpreting trajectories (Chapter~\ref{chap:dyno}), a robust tool for inferring linear trajectories (Chapter~\ref{chap:scorpius}), a novel way for studying gene regulation at a single-cell or single-sample level (Chapter~\ref{chap:bred}), and tool for quantifying the topological changes resulting from minor perturbations in a network (Chapter~\ref{chap:incgraph}).

We reflect on current challenges within the field, by providing guidelines for how to self-assess trajectory inference methods (Chapter~\ref{chap:selfassessment}) and, more generally, how to perform large-scale benchmarking experiments of computational methods (Chapter~\ref{chap:guidelines}).

Throughout this project I particularly enjoyed working in an open-source ecosystem, collaborating with scientists around the world to tackle unsolved problems, whilst learning best practices by inspecting the source code of exemplary software projects. 


\chapter*{Samenvatting}


Recente technologische ontwikkelingen hebben geleid tot significante ontwikkelingen in celbiologie die ons toelieten om de biomoleculaire inhoud van individuele cellen op een genoomwijde manier te observeren. Dankzij snelle ontwikkelingen in \textit{single-cell omics} technologieën wordt ons verstaan van cellulaire heterogeniteit en dynamische processen voortdurend gewijzigd. Inmiddels zijn er nieuwe spannende ontwikkelingen onderweg, zoals het \textit{Human Cell Atlas} project, die mikt om alle celtypes in het menselijke lichaam op biomoleculaire en spatiële wijze in kaart te brengen.

Gedurende dit doctoraatsproject heb ik deze inspirerende ontwikkelingen niet alleen van heel dichtbij kunnen volgen, maar ik heb ook wetenschappelijke bijdragen aan het veld kunnen leveren door nieuwe analytische programma’s extensief te onderzoeken, en door zelf verscheidene algoritmes te ontwikkelen. Meer specifiek heeft mijn onderzoek geleid tot de volgende bijdragen.

Een recurrent probleem van baanbrekende computationele algoritmes is dat gedurende hun ontwikkeling er helaas vaak onvoldoende datasets beschikbaar zijn om de performantie op een kwantitatieve wijze te kunnen beoordelen. We ontwikkelden een programma om individuele cellen te simuleren, dat flexibel en uitbreidbaar genoeg is om toekomstige computationele onderzoeksvelden te ondersteunen (Hoofdstuk~\ref{chap:dyngen}).

We voerden een comprehensieve studie van 45 \textit{trajectvoorspellende} methoden uit (Hoofdstuk~\ref{chap:dynbenchmark}). Trajectvoorspellende methoden is een nieuwe soort van computationele analyse waarbij het doel is om transities tussen verscheidene celtypes te identificeren en bestuderen. Terwijl deze klasse van methoden een van de snelst groeiende subdomeinen binnen \textit{single-cell omics} was, was er tot op dat moment nog geen kwantitatieve studie van de voor- en nadelen van dergelijke methoden uitgevoerd.

We ontwikkelde ook zelf vier nieuwe computer programma's, waaronder een uitgebreide toolkit om trajecten te voorspellen, visualiseren en interpreteren (Hoofdstuk~\ref{chap:dyno}), een robuust programma om lineaire trajecten te voorspellen (Hoofdstuk~\ref{chap:scorpius}), een nieuwe manier om gen regulatie op het niveau van een individuele cel of staal te bestuderen (Hoofdstuk~\ref{chap:bred}), en een programma om topologische veranderingen te kwantificeren ten gevolge van kleine veranderingen in een netwerk (Hoofdstuk~\ref{chap:incgraph}). 

We reflecteren op huidige uitdagingen binnen het veld, door richtlijnen aan te bieden over hoe je je eigen trajectvoorspellende methode kan evalueren (Hoofdstuk~\ref{chap:selfassessment}) en, in het algemeen, hoe je grootschalige beoordelingen van computationele methoden kan uitvoeren (Hoofdstuk~\ref{chap:guidelines}).

Tijdens het verloop van dit project heb ik bijzonder genoten van te werken in een \textit{open-source} omgeving, waarbij ik wereldwijd collaboraties kon aangaan met onderzoekers. Hierbij heb ik bijgeleerd door naar de broncode van voorbeeldige softwareprojecten te kunnen kijken.
